\section{Coherent moments and the limits of quadratic acquisition}
\label{sec:coherent-moments}

The logarithmic transfer preserves cancellations. Replacing its
integral by a quadratic moment loses information in two stages: the
oscillating reflected phase and the variation of the arithmetic arm.
We first compute the exact coherent moment, then quantify both losses.
The statements in this section are the E03 and B03 results
\cite{GNPRSMoment,GNPRSPhaseObstruction}.

\subsection{The exact local moment law}
Set \(M(X,c)=\int_{W_X}|V_{c/L_X}(x)|^2\,d\omega_X\).
\begin{theorem}[Coherent moment at \(c=1\)]\label{thm:coherent-moment}
For the complete field, with time fixed in each defining integral,
\begin{equation}\label{eq:coherent-moment}
 \lim_{T\to\infty}\frac1T\int_T^{2T}M(X,1)\dd X=\zeta(3/2).
\end{equation}
\end{theorem}
\begin{proof}
Introduce only for the height average
\(v(u)=V_{1/L_u}(u)\). Its auxiliary arm weights and phases are
\[
 w_n(u)=n^{-3/4}\ee^{(\log n)^2/(4L_u)}\le n^{-5/8},\quad
 \nu(u)=u/2+\pi/(8L_u),\quad
 \phi(u)=u(L_u-1)/2+\pi/8+7\pi/4.
\]
Extend weights by zero above \(N_u\). Each arm is \(O(u^{3/16})\).
By \eqref{eq:real-provider-error}, replacing \(|v|^2\) by the
square of the two arms costs \(O(u^{-1/16}+u^{-1/2})\).
The diagonal of the two arm squares is \(2\sum_{n\le N_u}w_n(u)^2\).
For every fixed \(n\), its summand tends uniformly on \([T,2T]\)
to \(2n^{-3/2}\); the dominating series \(2\sum n^{-5/4}\)
is summable. Thus this diagonal tends uniformly to \(2\zeta(3/2)\).

For every pair, integrate on its actual activation interval
\[
 I_{mn}=[T,2T]\cap[4\pi\max(m,n)^2,\infty).
\]
There \(w_mw_n\) decreases, and its supremum plus total variation
is at most \(2(mn)^{-5/8}\). A phase with monotone derivative of
fixed sign and modulus at least \(\alpha\) has every partial
integral bounded by \(3/\alpha\): integration by parts pays two
endpoints and the variation of the reciprocal derivative. Applying
this bound through the amplitude primitive pays the stated variation.

For same-arm pairs \(m\ne n\), eventually \(\nu'\ge1/4\) and
\(\nu''>0\). The cost is
\(O((mn)^{-5/8}/|\log(n/m)|)\). For \(m\le n/2\) use
\(|\log(n/m)|\ge\log2\); for \(n/2<m<n\) use
\(\log(n/m)\ge(n-m)/n\). Summing through \(N_{2T}\ll\sqrt T\)
gives an averaged error \(O(T^{-5/8}(1+\log T))\).

The cross-arm phases are \(\phi(u)-\nu(u)\log(mn)\), including
\(m=n\). On their activation intervals,
\[
 \frac1{8T}\le(\phi-\nu\log(mn))''\le\frac1{2T}
\]
eventually. The part where the first derivative has modulus at most
\(T^{-1/2}\) has length at most \(16\sqrt T\). Its complement
has at most two pieces, each costing \(3\sqrt T\). This gives
\(22\sqrt T\) for every partial integral, including a stationary
point at a cutoff endpoint. The complete averaged cross-arm cost is
\[
 O\left(T^{-1/2}\left(\sum_{n\le N_{2T}}n^{-5/8}\right)^2\right)
       =O(T^{-1/8}).
\]
We have proved \(T^{-1}\int_T^{2T}|v(u)|^2\dd u\to2\zeta(3/2)\).

It remains to restore the original clock and convolution. For
\(x\in W_X\), \(|1/L_X-1/L_x|\le4R_X/(XL_X^2)\).
At intermediate times \(1/2\le\lambda L_x\le2\), the common
weight bound is \(n^{-9/16}\); hence
\(S_\lambda=O(x^{7/32})\) and
\(\partial_\lambda S_\lambda=O(L_x^2x^{7/32})\).
The squared-arm comparison costs \(O(R_XX^{-9/16})\); paying the
full correction at both endpoint times adds \(O(X^{-1/32})\).
Thus it changes the local moment by a uniformly vanishing quantity.
No time-dependent field has been substituted without this payment.

Fubini now gives an auxiliary weight
\[
 W_T(x)=\int_T^{2T}K(x-X)\mathbf1_{|x-X|\le R_X}\dd X.
\]
It lies in \([0,1/2]\), and on
\([T+R_{2T},2T-R_{2T}]\) it is at least \(q_T\).
The remaining edge intervals have total length at most \(4R_{2T}\).
Since \(|v|^2=O(T^{3/8})\) there, replacing \(W_T\) by
\(\tfrac12\mathbf1_{[T,2T]}\) costs
\(O(T^{-29/8}+T^{-5/8}\log T)\) in the average.
The kernel mass \(1/2\) proves \eqref{eq:coherent-moment}.
\end{proof}

\subsection{An upper allowance that cannot succeed at every center}
Jensen's inequality applied with probability measure
\(d\omega_X/q_X\), together with the exact exterior, gives
\begin{equation}\label{eq:quadratic-allowance}
 J_\lambda(X)\le\mathcal B_J(X,c):=
 \frac{q_X}{2}\log\frac{M(X,c)}{q_X}
      -P_{c,\lambda}(X)+T_X+\frac{C_{\rm out}}{L_X}.
\end{equation}
Here \(M>0\), since an entire nonzero field cannot vanish on an
interval. The following failure concerns this upper allowance.

\begin{corollary}[Exact quadratic-Jensen obstruction]\label{cor:Jensen-stop}
There are arbitrarily large \(X\) such that
\(\mathcal B_J(X,1)\ge m_*\). The same centers fail if the exact
moment is replaced by any valid upper bound for it.
\end{corollary}
\begin{proof}
The elementary bounds needed are
\[
 m_*<\frac9{50},\qquad |P_{1,\lambda}|<\frac{23}{100},\qquad
 \zeta(3/2)>\frac{13}{5},\qquad \ee^{41/25}<\frac{26}{5}.
\]
For the first, use \(\sin(\pi/8)<77/200\) and
\(\cosh(\pi/8)>137/128\). For the second, the prime powers through
16 are \(2,3,4,5,7,8,9,11,13,16\); respective rational upper
budgets for their coefficients are
\(.14,.049,.012,.012,.005,.0014,.0013,.0014,.001,.00013\).
They sum to \(.22323\), and the remaining integer tail is at most
\(1/224<.0045\). These rational budgets follow by raising the
elementary fourth-root comparisons to the fourth power, as detailed
in \cite{GNPRSMoment}. Convexity of \(u^{-3/2}\) gives
\(\zeta(3/2)\ge1+1/(2\sqrt2)+13/(6\sqrt3)>177/68>13/5\).
Finally \(\ee^a\le1+a+a^2/[2(1-a/3)]\), \(0\le a<3\), follows
from \(k!\ge2\cdot3^{k-2}\); apply it at \(a=41/100\) and
raise to the fourth power.

If the allowance were below \(m_*\) eventually, dropping its
nonnegative errors would give
\[
 M(X,1)<q_X\exp\{(41/50)/q_X\}\longrightarrow
                  \tfrac12\ee^{41/25}<13/5.
\]
This eventual bound contradicts Theorem~\ref{thm:coherent-moment}.
Monotonicity of the logarithm proves the assertion for larger moments.
\end{proof}

\subsection{Why the ideal cosine correction is still insufficient}
For an independent uniform phase \(\vartheta\) and \(A\ne0\),
\[
 \mathbb E_\vartheta\log|2A\cos\vartheta|=\log|A|,
 \qquad \mathbb E_\vartheta|2A\cos\vartheta|^2=2|A|^2.
\]
The logarithmic identity follows from the sine/cosine doubling
formula with its integrable zero logarithms. If \(A\) also varies
under \(d\omega_X/q_X\), the product-average Jensen gap is
\begin{equation}\label{eq:phase-amplitude-gap}
 \frac{q_X}{2}\log2+
 q_X\left(\frac12\log\mathbb E_X|A|^2-\mathbb E_X\log|A|\right).
\end{equation}
The second term is nonnegative. This identity assumes an independent
phase average, not the actual spatial phase. Even granting its first
term for free does not repair the quadratic method.

\begin{theorem}[Phase-only correction obstruction]\label{thm:phase-stop}
Define the diagnostic
\(\mathcal B_{\rm ph}=\mathcal B_J-(q_X/2)\log2\).
There are arbitrarily large \(X\) with
\begin{equation}\label{eq:phase-stop}
 \mathcal B_{\rm ph}(X,1)\ge m_*+1/1000.
\end{equation}
No assertion that \(\mathcal B_{\rm ph}\) bounds \(J\) is needed
or made. The failure persists for any upper replacement of \(M\).
\end{theorem}
\begin{proof}
At \(c=1\), write \(P_0(X)=P_{1,0}(X)=\sum_p F_p((X/2)\log p)\),
where
\[
 F_p(\theta)=\sum_{k\ge1}
    \frac{p^{-11k/4}\cos(k\theta)}{k(1+p^{-4k})},\qquad
 b_p=\sum_{k\ge1}\frac{p^{-11k/4}}{k(1+p^{-4k})}.
\]
Each block has mean zero and lies in \([-b_p,b_p]\).
For a finite set of primes the height average of \(\exp(4\sum F_p)\)
equals the product circle average. Indeed the absolutely convergent
Fourier expansion has nonzero frequency
\(\tfrac12\sum j_p\log p\) for every nonconstant mode, by unique
factorization, and such a mode has vanishing dyadic mean.
The bound \(\sum b_p<\infty\) permits removal of the prime tail
uniformly. Convexity's chord bound and \(\cosh u\le\ee^{u^2/2}\)
therefore give
\begin{equation}\label{eq:prime-exponential-mean}
 \mathfrak E:=\lim_{T\to\infty}\frac1T\int_T^{2T}\ee^{4P_0(X)}\dd X
 \le\exp\left(8\sum_p b_p^2\right)<\ee^{27/125}.
\end{equation}
For completeness the last constant uses
\[
 b_2<77/500,\quad b_3<51/1000,\quad
 \sum_{p\ge5}b_p^2<\frac{2432}{7021125},\quad
 \sum_p b_p^2<\frac{27}{1000}.
\]
The first coefficients at 2 and 3 are below \(.14,.049\).
Writing \(r=p^{-11/4}\), the higher powers cost at most
\(r^2/[2(1-r)]\), with \(r_2<3/20,r_3<1/20\).
For \(p\ge5\), \(b_p<(80/79)p^{-11/4}\); sum its square over
integers using
\(\sum_{n\ge5}n^{-11/2}<1/6250+1/5625\).
This proves the displayed rational bounds.

Suppose \eqref{eq:phase-stop} fails eventually. Rearranging gives
\[
 M(X,1)<2q_X\exp\{2(m_*+1/1000+P_{1,1/L_X}(X))/q_X\}.
\]
The differentiated absolutely convergent prime series shows
\(P_{1,1/L_X}-P_0=O(1/L_X)\). Thus the right side is
\(\ee^{4(m_*+1/1000)}\ee^{4P_0(X)}+o(1)\), uniformly in \(X\).
Taking means and using \eqref{eq:prime-exponential-mean} forces
\[
 \zeta(3/2)<\ee^{47/50}
       \le\frac{26609}{10300}<\frac{13}{5},
\]
contrary to the preceding lower bound. The exponential bound uses
the same factorial inequality at \(47/50\).
\end{proof}

These obstructions locate information lost by the specified quadratic
expressions. They do not lower-bound the true logarithm or the Green
potential at the failing centers. In particular, no collision exclusion
or obstruction to all signed nonlinear estimates follows.

\subsection{Exact regularized phase and the fixed-prefix limit}
For any complex \(A\), put
\[
 B_\delta(A)=\frac{\delta+\sqrt{\delta^2+4|A|^2}}2,
 \qquad r_\delta(A)=\frac{\overline A^{,2}}{B_\delta(A)^2}.
\]
Since \(B_\delta-|A|^2/B_\delta=\delta\), direct expansion gives
\begin{equation}\label{eq:regularized-phase-factor}
 \tfrac12\log(|A+\ee^{\ii\phi}\overline A|^2+\delta^2)
 =\log B_\delta(A)+\Re\log(1+\ee^{\ii\phi}r_\delta(A)).
\end{equation}
Here \(|r_\delta|<1\), including \(A=0\), so the second logarithm
has its convergent power series. Its independent phase mean is zero.

For a fixed prefix \(A_{\lambda,M}\) and fixed \(\delta>0\),
\(|r_\delta|\le\rho<1\) and \(|\partial_x r_\delta|=O_{M,\delta}(1)\),
while \(\phi_\lambda'\ge L_X/4\) and
\(\phi_\lambda''=O(X^{-1})\). Integration by parts in the \(k\)-th
harmonic, including both endpoints, bounds its integral by
\[
 \frac{\rho^k(2\|K\|_\infty+\|K'\|_1)
       +k\rho^{k-1}q_X\|r_\delta'\|_\infty}{k\min\phi_\lambda'}
 +\frac{q_X\rho^k\|\phi_\lambda''\|_\infty}
       {k(\min\phi_\lambda')^2}.
\]
Summing with the logarithmic coefficient \(1/k\) gives
\(O_{M,\delta}(1/L_X)\). If
\(E_M=\int_{W_X}|V_\lambda-S_{\lambda,M}|^2\,d\omega_X\),
Lemma~\ref{lem:upper-log-transfer} consequently yields
\begin{align}\label{eq:prefix-transfer}
 J_\lambda(X)\le{}&\int_{W_X}\log B_\delta(A_{\lambda,M})\,d\omega_X
       -P_{c,\lambda}(X)
       +\frac{q_X}{2}\log\left(1+\frac{E_M}{q_X\delta^2}\right)
       \nonumber\\
 &+T_X+\frac{C_{\rm out}+C_{M,\delta}}{L_X}.
\end{align}
Restricting the moment proof to modes \(n>M\) gives exactly
\(T^{-1}\int_T^{2T}E_M(X,1)\dd X\to\sum_{n>M}n^{-3/2}\).
The tail is positive for each fixed prefix. This limits a vanishing
residual claim but does not decide the fixed defect margin for a
chosen prefix and regulator.

The fixed-prefix phase constants cannot be used for the full growing
cutoff. Already the relative phase at \(m=n=N_x\) has derivative
\((L_x-2\log N_x)/2+\pi\lambda/(8x)=O(x^{-1/2})\).
A large carrier derivative alone does not bound this derivative away
from zero. The next section chooses a different reference for which
the complete phase can be controlled.
